\chapter{Conclusão}
\label{chap:conclusion}

Este trabalho desenvolveu e avaliou uma infraestrutura de comunicação entre \textit{gateways} situados em redes privadas distintas, integrando o transporte QUIC aos mecanismos de conectividade do ICE. A solução partiu da arquitetura do Chownat e evoluiu para interfaces modulares que permitem combinar protocolos e caminhos de transporte, utilizando as bibliotecas \texttt{OpenSSL} e \texttt{libnice}. As aplicações mantêm suas conexões TCP locais, enquanto seus dados são transportados entre os \textit{gateways} por QUIC, diretamente ou por retransmissão TURN. A arquitetura do WebRTC serviu como referência para organizar as etapas de sinalização e estabelecimento da comunicação.

A avaliação funcional verificou a descoberta de endereços externos, incluindo portas diferentes das locais, o estabelecimento das conexões e a preservação dos dados nos dois sentidos. O teste prolongado demonstrou a continuidade da transferência durante a renovação de recursos TURN e sua remoção no encerramento, incluindo o tratamento do erro 438 após a expiração do \textit{nonce}. A avaliação em infraestrutura física também confirmou o estabelecimento QUIC sobre um caminho ICE direto e a transferência íntegra dos dados entre aplicações TCP locais. Esses resultados sustentam o atendimento dos objetivos de integração, modularidade e funcionamento do transporte nas configurações avaliadas. A comparação conceitual evidenciou que o caminho direto dispensa a retransmissão contínua por um servidor, enquanto o uso de TURN mantém essa dependência, preservando a proteção QUIC entre os participantes.

As conclusões restringem-se às topologias, cargas e condições exercitadas. Os testes com TURN obrigatório verificaram a retransmissão, mas não a seleção automática desse caminho após uma tentativa direta malsucedida. Não foram realizadas medições comparativas de latência ou vazão, nem avaliações sistemáticas com perdas induzidas, mobilidade ou dispositivos IoT com recursos limitados. A contribuição concentra-se, portanto, na infraestrutura de transporte entre os \textit{gateways}, sem abranger a validação de um sistema IoT completo.

Como continuidade, destacam-se a ampliação dos cenários de NAT e a avaliação de desempenho e consumo de recursos. Outra linha de investigação consiste em desenvolver mecanismos de \textit{hole punching} baseados na migração de conexão do QUIC que preservem as restrições do protocolo e a proteção contra ataques de amplificação, enfrentando as limitações identificadas nas propostas analisadas. Também se propõe explorar um protocolo próprio inspirado no QUIC, voltado ao transporte genérico entre \textit{gateways}, com entrega confiável, controle de fluxo e congestionamento e proteção criptográfica. Essa alternativa exigiria especificar, implementar e avaliar as garantias oferecidas e sua integração aos mecanismos de travessia de NAT.